\documentclass[journal,onecolumn]{IEEEtran}

% ── Encabezado ──────────────────────────────────────────────
\usepackage{fancyhdr}
\pagestyle{fancy}
\fancyhf{}
\renewcommand{\headrulewidth}{0.4pt}
\fancyhead[L]{\small ERS/SRS Final}
\fancyfoot[C]{\thepage}

% ============================================================
% IDIOMA Y CODIFICACIÓN  (mover ANTES de csquotes)
% ============================================================
\usepackage[utf8]{inputenc}
\usepackage[T1]{fontenc}
\usepackage[spanish,es-tabla]{babel}

% ============================================================
% CITAS Y TEXTO
% ============================================================
\usepackage{csquotes}

% ============================================================
% ESPACIADO Y LISTAS
% ============================================================
\usepackage{parskip}
\setlength{\parskip}{4pt}
\usepackage{enumitem}

% ============================================================
% IDIOMA Y CODIFICACIÓN
% ============================================================
\usepackage[spanish,es-tabla]{babel}
\usepackage[utf8]{inputenc}
\usepackage[T1]{fontenc}

% ============================================================
% PAPEL Y GEOMETRÍA
% ============================================================
\usepackage{geometry}
\geometry{a4paper, top=2.5cm, bottom=2.5cm, left=2.5cm, right=2.5cm}

% ============================================================
% IMÁGENES
% ============================================================
\usepackage{graphicx}
\graphicspath{{Figuras/}}

% ============================================================
% TABLAS
% ============================================================
\usepackage{booktabs}
\usepackage{tabularx}
\usepackage{longtable}

% ============================================================
% MATEMÁTICAS
% ============================================================
\usepackage{amsmath}
\usepackage{amssymb}

% ============================================================
% CÓDIGO FUENTE
% ============================================================
\usepackage{listings}

% ============================================================
% FORMATO ACADÉMICO
% ============================================================
\usepackage{microtype}
\usepackage{float}

% ============================================================
% CAPTIONS (TÍTULOS DE FIGURAS/TABLAS)
% ============================================================
\usepackage{caption}
\captionsetup{font=small, labelfont=bf}

% ============================================================
% REFERENCIAS / HIPERVÍNCULOS
% ============================================================
\usepackage[hidelinks]{hyperref}

% ============================================================
% NUMERACIÓN Y NOMBRES EN ESPAÑOL
% ============================================================
\renewcommand{\thetable}{\arabic{table}}
\def\tablename{Tabla}

\renewcommand{\thesection}{\arabic{section}}
\renewcommand{\thesubsection}{\thesection.\arabic{subsection}}
\renewcommand{\thesubsubsection}{\thesubsection.\arabic{subsubsection}}

% ============================================================
% ESTILO DE TÍTULOS DE SECCIÓN
% ============================================================
\usepackage{titlesec}
\titleformat{\section}{\normalfont\Large\bfseries}{\thesection}{1em}{}
\titleformat{\subsection}{\normalfont\large\bfseries}{\thesubsection}{1em}{}
\titleformat{\subsubsection}{\normalfont\normalsize\bfseries}{\thesubsubsection}{1em}{}

% ============================================================
% CADA SECCIÓN PRINCIPAL (\section) INICIA EN PÁGINA NUEVA
% ============================================================
\usepackage{etoolbox}
\pretocmd{\section}{\clearpage}{}{}

% ============================================================
% BIBLIOGRAFÍA (BIBLATEX + BIBER)
% ============================================================
\usepackage[
backend=biber,
style=ieee,
sorting=none
]{biblatex}
\addbibresource{referencias_estructura.bib}

% ============================================================
% DOCUMENTO
% ============================================================
\begin{document}

% ============================================================
% SECCIÓN 3: ERS/SRS FINAL DEL SISTEMA MUNDIPETS
% Responsable: GUTIÉRREZ ORTEGA, Génesis Adriana
% ============================================================

\section{ERS/SRS Final del Sistema MundiPets}
\label{sec:ers-final}

\subsection{Estructura conforme a ISO/IEC/IEEE 29148:2018}
\label{sec:estructura-29148}

El ERS final de MundiPets se organizó siguiendo la estructura de contenido que
ISO/IEC/IEEE 29148:2018 establece para una Especificación de Requisitos de
Software \cite{ieee29148}. La Tabla~\ref{tab:correspondencia-29148} muestra la
correspondencia explícita entre las cláusulas de la norma y las secciones
efectivamente desarrolladas en el documento del equipo; no se transcribe aquí
el contenido de cada sección, que se encuentra íntegro en el ERS consolidado
del repositorio, sino que se documenta dónde localizarlo y bajo qué cláusula
normativa queda encuadrado.

\begin{table}[H]
	\centering
	\caption{Correspondencia entre las cláusulas de ISO/IEC/IEEE 29148:2018 y las secciones del ERS de MundiPets.}
	\label{tab:correspondencia-29148}
	\footnotesize
	\renewcommand{\arraystretch}{1.3}
	\begin{tabularx}{\textwidth}{
		|>{\raggedright\arraybackslash}p{3.6cm}
		|>{\raggedright\arraybackslash}X|}
		\hline
		\textbf{Cláusula ISO/IEC/IEEE 29148:2018} & \textbf{Sección del ERS de MundiPets} \\
		\hline
		1. Introducción & Sección 1 --- Introducción \\
		\hline
		\quad 1.1 Propósito & 1.1 Propósito del documento \\
		\hline
		\quad 1.2 Alcance & 1.2 Alcance del producto \\
		\hline
		\quad 1.3 Definiciones, acrónimos y abreviaturas & 1.3 Definiciones, acrónimos y abreviaturas (glosario) \\
		\hline
		\quad 1.4 Referencias & 1.4 Referencias \\
		\hline
		\quad 1.5 Visión general del documento & 1.5 Visión general del documento \\
		\hline
		2. Descripción general & Sección 2 --- Descripción general \\
		\hline
		\quad 2.1 Perspectiva del producto & 2.1 Perspectiva del producto y límite del sistema (incluye el entorno operativo) \\
		\hline
		\quad 2.2 Funciones del producto & 2.2 Funciones del producto \\
		\hline
		\quad 2.3 Características de los usuarios & 2.3 Clases y características de usuarios; 2.4 Mapa de stakeholders (matriz poder/interés y conflictos potenciales) \\
		\hline
		\quad 2.4 Restricciones generales & 2.5 Entorno operativo (aspectos no cubiertos en 2.1) \\
		\hline
		\quad 2.5 Suposiciones y dependencias & 2.6 Suposiciones y dependencias \\
		\hline
		3. Requisitos específicos & Sección 3 --- Requisitos específicos completos \\
		\hline
		\quad 3.1 Requisitos de interfaces externas & 3.1 Interfaces externas (usuario, hardware, software) \\
		\hline
		\quad 3.2 Requisitos funcionales & 3.2 Requisitos funcionales (RF-01 a RF-27) \\
		\hline
		\quad 3.3 Requisitos de rendimiento & Incorporados dentro de 3.3 Requisitos no funcionales (RNF-02, RNF-03, RNF-12) \\
		\hline
		\quad 3.4 Restricciones de diseño & 3.5 Restricciones de diseño (RD-01 a RD-09) \\
		\hline
		\quad 3.5 Atributos del sistema software (calidad) & 3.3 Requisitos no funcionales (RNF-01 a RNF-16) \\
		\hline
		\quad 3.6 Otros requisitos (legales, de negocio) & 3.4 Requisitos legales \\
		\hline
		4. Información de apoyo & 3.6 Historias de usuario y criterios de aceptación (verificabilidad); Sección 4 --- Modelado del sistema con UML; Sección 5 --- Priorización y trazabilidad extendida (matriz de trazabilidad) \\
		\hline
	\end{tabularx}
\end{table}

La correspondencia es completa: ninguna cláusula de contenido que la norma
exige para un SRS quedó sin sección equivalente en el ERS de MundiPets. La
distinción entre requisitos de rendimiento (cláusula 3.3) y atributos de
calidad del software (cláusula 3.5) que la norma trata como categorías separadas se resolvió integrando ambas dentro de la sección única de
Requisitos no funcionales del documento del equipo, clasificando internamente
cada RNF según corresponda a rendimiento, seguridad, usabilidad, mantenibilidad
u otro atributo de calidad; esta es una decisión de organización del equipo,
no una omisión, y se declara aquí de forma explícita por la misma razón que
exige la verificabilidad de un requisito individual: que la trazabilidad entre
lo declarado y lo exigido sea comprobable por cualquier lector, incluido el
tribunal.

El ERS del equipo contiene, además de las cláusulas normativas anteriores, dos
bloques de contenido que exceden el mínimo que exige ISO/IEC/IEEE 29148:2018:
la Sección 6, Producto Mínimo Viable (alcance del MVP, arquitectura de
despliegue y evidencia de funcionamiento), y la Sección 7, Componente empírico
del proyecto (protocolo de validación experimental, preguntas de investigación
en formato PICOC, hipótesis, variables, participantes, instrumentos, plan de
análisis estadístico y registro previo del protocolo). Ninguno de estos dos
bloques corresponde a una cláusula de la norma y, por lo tanto, no se fuerza su
inclusión en la Tabla~\ref{tab:correspondencia-29148}: son contenido adicional
que el equipo decidió mantener como valor agregado del proyecto integrador,
más allá de la estructura mínima que un SRS conforme a 29148 exige.

\printbibliography[heading=none]

\end{document}
